A language model can interpret an acronym in a sentence by generating an expansion or by choosing from an inventory. We study this distinction for Hebrew using five systems: a fine-tuned DictaBERT candidate ranker, and generation and selection configurations of Qwen and Gemini. On a frozen, type-disjoint held-out test set of natural Hebrew acronym occurrences, free generation performs poorly for both LLMs, largely failing to produce the inventory's exact wording; restricting the same models to select from a fixed candidate list closes most of this gap. The fine-tuned encoder clearly outperforms unconstrained LLM generation but does not surpass the strongest LLM under selection. These results suggest that, for this task, a general-purpose LLM does not already solve Hebrew acronym disambiguation from context alone, but comes close once given the same candidate scaffold a task-trained system relies on. We describe the systems, data construction and known limitations that inform this comparison.
